\documentclass[../../main.tex]{subfiles}
\begin{document}

\begin{center}
  {\fontsize{48pt}{58pt}\selectfont \textbf{1961}\normalsize\textbf{年度}}
\end{center}

\begin{table}[htb]
  \centering
  \caption{出題テーマと難易度}
  \begin{tabular}{|c|p{9cm}|c|}
    \hline
    \textbf{問題番号} & \textbf{テーマ} & \textbf{難易度} \\
    \hline
    第1問 & 3本の半直線上を動く点と三角形 & \\
    \hline
    第2問 & 四次関数の値から求める導関数 & \\
    \hline
    第3問 & 半球に外接する円錐の表面積の最小化 & \\
    \hline
    第4問 & 三角形内の線分と面積比 & \\
    \hline
    第5問 & 媒介変数曲線と座標軸に囲まれる面積 & \\
    \hline
    第6問 & 三角関数の係数と最小値 & \\
    \hline
  \end{tabular}
\end{table}

\begin{table}[htb]
  \centering
  \caption{解答}
  \renewcommand{\arraystretch}{1.4}
  \begin{tabular}{|c|c|p{9cm}|}
    \hline
    \textbf{問題番号} & \textbf{小問} & \textbf{答え} \\
    \hline
    第1問 & (1) & $t=\dfrac{1+\sqrt{5}}{2}$ \\
    \cline{2-3}
    & (2) & $t=0,1,2$ \\
    \hline
    第2問 & -- & $f'(0)=2.715$ \\
    \hline
    第3問 & -- & 高さ$2a$ \\
    \hline
    第4問 & -- & 面積比は$1:7$ \\
    \hline
    第5問 & -- & $\dfrac{9}{2}$ \\
    \hline
    第6問 & (1) & $(a,b,c)=(5,5,\sqrt{2})$ \\
    \cline{2-3}
    & (2) & $x=\dfrac{5\pi}{4}$で最小値$-4\sqrt{2}$ \\
    \hline
  \end{tabular}
\end{table}

\paragraph{試験形態}
6問．試験時間は要確認．

\end{document}
